% Generated draft: requires mathematical and source review.
\documentclass[11pt]{article}
\usepackage[margin=25mm]{geometry}
\usepackage{fontspec}
\setmainfont{Latin Modern Roman}
\usepackage{amsmath,amssymb,mathtools,unicode-math}
\setmathfont{Latin Modern Math}
\usepackage{enumitem,needspace,xurl,hyperref}
\setlength{\parindent}{0pt}
\setlength{\parskip}{5pt}
\setlength{\emergencystretch}{3em}
\newcommand{\N}{\mathbb{N}}
\newcommand{\Z}{\mathbb{Z}}
\newcommand{\Q}{\mathbb{Q}}
\newcommand{\R}{\mathbb{R}}
\newcommand{\C}{\mathbb{C}}
\newcommand{\F}{\mathbb{F}}
\newcommand{\A}{\mathbb{A}}
\newcommand{\PP}{\mathbb{P}}
\newcommand{\E}{\mathbb{E}}
\newcommand{\eps}{\varepsilon}
\newcommand{\dd}{\,\mathrm{d}}
\begin{document}
\section*{UnsolvedMath ID: 30001093}
\textbf{Source problem code:} catalogue entry
\subsection{Definitions and mathematical statement}
Let $G_0,G_1$ admit finite CW models $K_i=K(G_i,1)$, so $\widetilde K_i$ is contractible and $\pi_1(K_i)=G_i$. Assume $b_j^{(2)}(G_i)=0$ in every degree. These Betti numbers are the group-von-Neumann dimensions of the reduced homology of the $\ell^2$ cellular chain complex of $\widetilde K_i$. To define finite real $L^2$-torsion, one also needs determinant-class conditions: the restricted cellular differentials must have finite logarithmic Fuglede--Kadison determinants. The resulting torsion is an alternating sum of these logarithmic determinants. Vanishing $L^2$-Betti numbers alone does not state this additional hypothesis. Moreover, writing the torsion as a group invariant $\rho^{(2)}(G_i)$ requires the appropriate independence from the chosen finite classifying model; the extract does not give the determinant/homotopy-invariance assumptions.\par Subject to these definitional qualifications, the source asks two alternative invariance questions: if $G_0$ and $G_1$ are quasi-isometric, must \[\rho^{(2)}(G_0)=0\quad\Longleftrightarrow\quad\rho^{(2)}(G_1)=0?\] Does the same equivalence hold if the groups are orbit equivalent? Quasi-isometry is understood for word metrics on these finitely generated groups: a coarsely surjective map distorts distances by bounded multiplicative and additive constants. The standard measured meaning of orbit equivalence is existence of free ergodic probability-measure-preserving actions on standard probability spaces and a measure-space isomorphism carrying orbits to orbits almost everywhere. The source does not spell out that action convention. Either equivalence is a separate proposed sufficient hypothesis; both are not required simultaneously. The claim concerns simultaneous vanishing, not equality of the numerical torsion values. The 2026-08-21 triage labels the record partially solved and reports Farrell--Jones closure and later torsion results in important cases, including uniform measure-equivalence invariance under determinant hypotheses, while leaving the full quasi-isometry/orbit-equivalence question unresolved. The stronger uniform-measure-equivalence and determinant assumptions in that progress report are not silently omitted or identified with arbitrary orbit equivalence.
\subsection{Short English statement}
Ask whether quasi-isometry, or separately orbit equivalence, preserves vanishing of $L^2$-torsion for groups with finite classifying spaces and vanishing $L^2$-Betti numbers. Determinant-class and model-independence conditions are needed to make the displayed group torsions well defined.
\subsection{Sources}
\begin{itemize}\item Local UnsolvedMath registry, permanent ID 30001093, source problem catalogue entry, \emph{Invariance of $L^2$-Torsion Under Large-Scale Equivalence}, proposed year 2008; no individual proposer supplied.\item Supplied citation: \emph{Von Neumann Algebras and Ergodic Theory of Group Actions} (2009), Oberwolfach Reports, DOI 10.4171/owr/2008/49, pp.\ 2763--2814: \url{https://doi.org/10.4171/owr/2008/49}. Supplied literature label: original OWR report.\item Supplied literature label: L\"uck--Sauer--Wegner, L2-torsion and uniform measure equivalence, \url{https://arxiv.org/abs/0903.2925}.\item Literature assessment date: 2026-08-21. The status and progress claims are supplied dated curation, not independently verified here.\end{itemize}
\end{document}
